% ======================================================================
% SECTION 1: INTRODUCTION
% Five thematic subsections: FDA RTCT, GBM clinical context, current
% robotic neurosurgery baseline, on-premises LLM thesis, and the
% transition to the 60-second resection.
% Bracketed prompts for future Claude Code Opus 4.7 1M Max processing.
% Do not process the bracketed prompts here.
% ======================================================================

\subsection{FDA Real-Time Clinical Trials Announcement}
\label{sec:intro-fda}

[INTRODUCTION SUBSECTION 1 PROMPT - PROCESS DURING THE PAPER POPULATION PASS]

[Open with the FDA's 28 April 2026 announcement that real-time clinical
trials are entering the proof-of-concept phase. Read the press release
text and key quotes from the prior paper's Introduction at
kevinkawchak/physical-ai-oncology-trials/new-trial/national-24-7-trial/paper/full-paper/final-paper/main.tex
and reuse only the FDA framing (do not copy any prose verbatim).
Anchor the framing to fda2026realtime in references.bib. Position the
60-second glioblastoma robotic surgery as a natural extension of the
RTCT vision toward the surgical theater, not a critique of it. Note
that the FDA's two proof-of-concept trials (AstraZeneca's TRAVERSE in
mantle cell lymphoma and Amgen's STREAM-SCLC in small cell lung
carcinoma) cover pharmacology only; the 60-second glioblastoma
resection extends real-time signal sharing into the surgical theater
through the on-premises LLM thesis. Mention the FDA Commissioner Marty
Makary and Chief AI Officer Jeremy Walsh quotes only by attribution
(do not reproduce the full quotes). Close the subsection by stating
that the United States must remain Number 1 in the world regarding
patient safety, efficacy, and speed benefits in oncological robotic
surgeries in clinical trials; the on-prem LLM thesis is one route to
that goal.]

\subsection{Glioblastoma Clinical Context}
\label{sec:intro-gbm}

[INTRODUCTION SUBSECTION 2 PROMPT - PROCESS DURING THE PAPER POPULATION PASS]

[Summarize the IDH-wildtype glioblastoma multiforme clinical context
relevant to the 60-second resection target: median survival 12 to 15
months under standard-of-care radiotherapy plus temozolomide (cite
Stupp2005Glioblastoma), the published extent-of-resection threshold
for newly diagnosed glioblastoma (cite Sanai2011GBMResection), and
5-aminolevulinic acid fluorescence-guided resection (cite
Stummer20055ALA). Read the file
kevinkawchak/physical-ai-oncology-trials/competitions/instructions/one_minute_variant/glioblastoma_context_1min.md
in full. Reproduce the patient profile (synthetic PAT-GBM-0001,
4.2 cm right frontal IDH-wildtype glioblastoma) and the maximal-safe
gross-total resection target. Highlight that current operative
duration ranges between 4 and 8 hours; the 60-second target is a 240x
to 480x compression that is currently not feasible with any commercial
robot. The contrast frames the rest of the paper as a checking step
the industry has not yet taken. Also reference the author's prior
glioblastoma and clinical trial work (cite kawchak_2025_17774560,
kawchak_2025_15549831, kawchak_2025_17614396, and kawchak_2026_19994945)
as the progression toward this 60-second computational target.]

\subsection{Current Robotic Neurosurgery State of the Art}
\label{sec:intro-baseline}

[INTRODUCTION SUBSECTION 3 PROMPT - PROCESS DURING THE PAPER POPULATION PASS]

[Read
kevinkawchak/physical-ai-oncology-trials/competitions/instructions/one_minute_variant/robot_specification_neurospeed.md
in full. Reproduce the side by side specification comparison between
the current state of the art Medtronic ROSA ONE Brain v3.0 and the
hypothetical 2030 Medtronic NeuroSpeed 1.0 across at least the
following parameters: tissue removal rate (mm cubed per second per
arm), end effector velocity (50 mm per second vs 1,000 mm per second;
20x), end effector acceleration (200 mm per second squared vs 10,000
mm per second squared; 50x), joint angular velocity, E-stop latency
(50 ms vs 5 ms; 10x), positioning accuracy at speed (0.5 mm RMS vs
0.1 mm RMS; 5x), force resolution (0.01 N vs 0.001 N; 10x), and duty
cycle. Render the comparison as a table using the column type
{>{\raggedright\arraybackslash}p{Xcm}} for each column to avoid
large gaps between words. Cite rosa-one-brain, Lefranc2014ROSAStereotaxy,
and Kalakoti2019RoboticsReview for the baseline. State explicitly
that the current state of the art is short of the 1-minute requirement
on every key parameter by 5 to 200 times. Note the liquid nitrogen
cooling design and 5-minute peak duty cycle of NeuroSpeed 1.0 that
close the gaps. Reference IEC 80601-2-77 (per-arm 5.0 N tip and 1.0 N
lateral force limits, 5 ms E-stop budget; cite iec-80601-2-77) and
IEC 62304 (safety-critical software lifecycle at the 1 kHz heartbeat
layer; cite iec-62304) as the regulatory floor.]

\subsection{On-Premises LLM Thesis}
\label{sec:intro-thesis}

[INTRODUCTION SUBSECTION 4 PROMPT - PROCESS DURING THE PAPER POPULATION PASS]

[State the thesis verbatim: on-premises repository based LLMs provide
commands to standard oncology surgical robots based on real-time
sensor data and controlled via x, y, z coordinates to administer
patient treatment, and this workflow minimizes single robot error
potential. Read kevinkawchak/robotic-surgeries/2030-gbm-1min/README.md
in full and reproduce the 4-arm coordination heartbeat ASCII snapshot
(1 kHz broadcast, 32-byte frame, 1 ms deadline, cumulative ee_force
across 4 arms less than or equal to 12 N, per-arm tip force less than
or equal to 5.0 N, E-stop 5 ms, heartbeat watchdog 3 ms). Map each
arm to its tool assignment: arm 1 hybrid ultrasonic plus waterjet
plus pulsed plasma; arm 2 bipolar coagulation plus irrigation; arm 3
suction plus tissue collection; arm 4 0.5 T iMRI plus 5-ALA
fluorescence plus ultrasound. Cite claude-opus-47 and claude-code for
the 1M token context window that makes on-premises orchestration
tractable. Cite ollama and vllm for the on-premises inference
backends compatible with the comparison agent. Close the subsection
with the practical insight: the on-prem LLM is a Software as a
Medical Device candidate under the FDA SaMD framework (cite fda-samd).
Highlight the safety advantage of single robot error minimization
through 4-arm cross-checking on the deterministic 1 kHz heartbeat
bus.]

\subsection{Transition to the 60-Second Resection}
\label{sec:intro-transition}

[INTRODUCTION SUBSECTION 5 PROMPT - PROCESS DURING THE PAPER POPULATION PASS]

[Introduce the 16-iteration deterministic sweep across noise (0.01 to
0.05 mm sensor noise sigma), gain (0.8 to 1.2 force feedback gain),
inverse kinematics tolerance (1e-6 to 1e-3), and heartbeat jitter (0
to 50 microseconds), read from
kevinkawchak/physical-ai-oncology-trials/competitions/instructions/one_minute_variant/commit_04_iterations_1min.md.
State the deterministic seed 20260510. State the per-iteration L1
50 ms, L2 1 s, L3 per-phase, and events Parquet aggregates emitted
to kevinkawchak/robotic-surgeries/2030-gbm-1min/outputs/iterations/.
Introduce the LLM tournament at
kevinkawchak/robotic-surgeries/2030-gbm-1min/outputs/comparison/
(default robot vs robot, 6 rounds) and the mixed tournament at
kevinkawchak/robotic-surgeries/2030-gbm-1min/outputs/comparison_robot_vs_human/
(2 robot plus 2 human, 6 rounds). Foreshadow the headline result:
robot mean composite score 88.53 vs human mean composite 70.35; robot
wins all 4 robot-vs-human pairings with confidence 0.955 to 1.000;
the structural-time-dimension caveat (1-minute robot vs 1-hour human)
is preserved in every rationale. Cite kawchak_2026_20113157 and
repo-robotic-surgeries for the artifact set. Close the subsection
with the on-screen Table of Contents pointer (the Table of Contents
follows on a fresh page).]
